\documentclass[11pt]{article}
\input{../common/preamble.tex}
\hypersetup{pdftitle={Optimal simplex products for projection-body volume},
  pdfauthor={Research manuscript prepared with ChatGPT for Vladimir Reshetnikov}}
\title{Optimal simplex products for projection-body volume\\[0.4em]
  \large Exact extremizers, periodicity, and a sharp isolation gap}
\author{Research manuscript prepared for the ProveIt project}
\date{8 October 2026}
\begin{document}
\maketitle
\begin{abstract}
We determine the maximum of normalized projection-body volume over
all affine images of Cartesian products of simplices in every
dimension. The unique optimal exponential block size is thirteen.
In every total dimension, only two balanced products need be compared.
From dimension one hundred onward, the maximizing product is given
by an exact period-thirteen rule, using thirteen-dimensional factors
and at most eight fourteen-dimensional or four twelve-dimensional
factors. We prove uniqueness of the maximizing dimension multiset and
a sharp, dimension-independent multiplicative gap to every other
product. The optimum satisfies an exact recurrence and has an explicit
periodic asymptotic expansion relative to the simplex value, with
optimal exponential base approximately $1.03372825553$.
All required finite inequalities are certified by integer arithmetic,
and an independent exact dynamic program checks the formulas through
dimension three hundred. The results refine a recent product
construction; they make no claim about the unrestricted maximum over
all convex bodies. The geometric reduction and every optimization
argument are proved in full.
\end{abstract}
\begingroup\small\tableofcontents\endgroup
\clearpage
\section{Introduction and scope}
\label{sec:projection-optimization}

For an $n$-dimensional convex body $K$, write
\[
 R_n(K)=\frac{|\Pi K|}{|K|^{n-1}},
 \qquad
 s_d=\frac{(d+1)d^d}{d!}.
\]
The projection body is defined by
$h_{\Pi K}(u)=|\operatorname{proj}_{u^\perp}K|_{n-1}$ for unit $u$.
We consider the class of invertible affine images of Cartesian products
of nondegenerate simplices. The factor dimensions can vary, and their
number is unrestricted. The results in this section optimize over that
class; they do not determine the maximum over all convex bodies.

The source manuscript \cite{OpenAIProjection} proves the two identities
\begin{equation}\label{eq:proj-inputs}
 R_d(T_d)=s_d,
 \qquad
 R_{r+s}(A\times B)=R_r(A)R_s(B),
\end{equation}
where $T_d$ is any nondegenerate $d$-simplex. Section~\ref{sec:proj-geometry} rederives them for completeness. The source uses repeated
ten-dimensional blocks to obtain an exponential improvement over a
single simplex. We optimize that construction exactly, identify its
unique optimal exponential block size, and give a sharp isolation gap
between optimal and nonoptimal products. The product counterexample
itself is not new: the source explicitly cites the earlier
counterexamples of Feng, Hu, Liu, and Xu \cite{FengReverse} to the
unrestricted simplex-maximum conjecture of Brannen~\cite{Brannen1996}.

Let
\begin{equation}\label{eq:proj-Pn}
 P_n=\max\left\{\prod_{i=1}^k s_{d_i}:
  k\ge1,\ d_i\ge1,\ \sum_{i=1}^k d_i=n\right\}.
\end{equation}
In dimension one the natural convention is $\Pi I=[-1,1]$ and $s_1=2$.
The facet proof of \eqref{eq:proj-inputs} extends to an interval factor
by assigning volume one to a point. Alternatively, interval factors can
be omitted throughout: whenever $n\ge2$, merging a dimension-one factor
with another factor strictly increases the objective, as follows from
the ratio formula below. Thus both conventions give the same $P_n$ for
$n\ge2$.

\section{Exact optimization}

Define
\begin{equation}\label{eq:proj-constants}
 \rho=s_{13}^{1/13},\quad L=\log\rho,\quad
 \alpha=12L-\log s_{12},\quad
 \beta=14L-\log s_{14}.
\end{equation}
The constants $\alpha$ and $\beta$ are positive. Numerically,
\[
 \rho\simeq2.80996473256,\qquad
 \alpha\simeq0.001448531660,\qquad
 \beta\simeq0.000775423091.
\]
The decimals are descriptive; all inequalities used in the proof have
exact rational certificates.

\begin{theorem}[Exact eventual optimizer]\label{thm:proj-optimum}
Write $n\equiv r\pmod {13}$, with $0\le r\le12$.
For every $n\ge100$,
\begin{equation}\label{eq:proj-exact}
 P_n=
 \begin{cases}
 s_{13}^{(n-14r)/13}s_{14}^{r},&0\le r\le8,\\[2mm]
 s_{12}^{13-r}s_{13}^{(n-12(13-r))/13},&9\le r\le12.
 \end{cases}
\end{equation}
Every maximizing factorization has exactly the indicated dimensions,
up to ordering. Consequently the maximizing bodies in the stated class
are uniquely specified up to invertible affine maps.

The first formula is valid as soon as $n\ge14r$; the second is valid as
soon as $n\ge12(13-r)$. The uniform threshold $100$ is the smallest
threshold for which these residue formulas hold in every dimension.
Moreover,
\begin{equation}\label{eq:proj-rate}
 \sup_{d\ge1}s_d^{1/d}=s_{13}^{1/13}=\rho,
\end{equation}
and equality in \eqref{eq:proj-rate} occurs only at $d=13$.
\end{theorem}

For all dimensions, not merely those beyond the periodic threshold, a
two-candidate formula suffices. Given $1\le k\le n$, put
\[
 a=\lfloor n/k\rfloor,\qquad b=n-ka,
 \qquad B(n,k)=s_a^{k-b}s_{a+1}^{b}.
\]

\begin{theorem}[Two candidates in every dimension]\label{thm:proj-two}
For every positive integer $n$,
\[
 P_n=\max\{B(n,k):
 k\in\{\max(1,\lfloor n/13\rfloor),\lceil n/13\rceil\}\}.
\]
Repeated candidates are counted once. In particular, computing the
optimal value requires comparing only two explicit rational numbers.
For $n\le19$ the simplex is optimal; the first strict product
improvement occurs at $n=20$, with dimensions $(10,10)$.
\end{theorem}

\section{Geometric reduction}\label{sec:proj-geometry}

We prove the geometric identities used in the optimization. They are
classical or explicitly proved in the source manuscript
\cite{OpenAIProjection}; the proofs here are included to make the
present argument independent of any unverified geometric assertion.
For a compact convex set $C$, its support function is
$h_C(u)=\max_{x\in C}\langle u,x\rangle$. Support functions determine
compact convex sets, and Minkowski addition adds their support
functions.

\begin{lemma}[Facet formula]\label{lem:proj-facets}
Let $K\subset\R^n$ be a full-dimensional polytope. For each facet $F$,
write $a_F$ for its $(n-1)$-volume and $\nu_F$ for its outward unit
normal. Then
\[
 h_{\Pi K}(u)=\frac12\sum_F a_F|\langle\nu_F,u\rangle|,
 \qquad
 \Pi K=\sum_F[-a_F\nu_F/2,a_F\nu_F/2].
\]
For $n=1$ a facet is an endpoint and has volume one.
\end{lemma}

\begin{proof}
Take a unit vector $u$ and project onto $u^\perp$. The absolute
Jacobian of projection from the hyperplane of $F$ is
$|\langle\nu_F,u\rangle|$. Apart from the projections of lower
dimensional faces and projection-boundary points, which have measure
zero, a line parallel to $u$ through a point of the shadow enters and
leaves $K$ through exactly two facet interiors. Summing projected facet
areas therefore counts the shadow twice. This proves the first
identity for unit $u$; homogeneity proves it for arbitrary $u$.
The support function of $[-v/2,v/2]$ is
$|\langle v,u\rangle|/2$, proving the second identity. In dimension
one the two endpoints give $h_{\Pi K}(u)=|u|$, as required.
\end{proof}

\begin{proposition}[Products and affine maps]\label{prop:proj-product}
Let $A\subset\R^r$ and $B\subset\R^s$ be full-dimensional polytopes,
with $r,s\ge1$. Then
\[
 \Pi(A\times B)=(|B|\Pi A)\times(|A|\Pi B),
 \qquad R_{r+s}(A\times B)=R_r(A)R_s(B).
\]
For every invertible linear map $U$ in dimension $n$,
\[
 \Pi(UK)=|\det U|U^{-t}\Pi K,
 \qquad R_n(UK)=R_n(K).
\]
Translations also preserve $R_n$.
\end{proposition}

\begin{proof}
A facet of $A\times B$ is either $F\times B$, with $F$ a facet of
$A$, or $A\times G$, with $G$ a facet of $B$. Indeed, an exposed face
is a product of exposed faces, and codimensions add. Their area-normal
vectors are respectively $(|B|a_F\nu_F,0)$ and
$(0,|A|a_G\nu_G)$. The facet formula gives
\[
 h_{\Pi(A\times B)}(u,v)
       =|B|h_{\Pi A}(u)+|A|h_{\Pi B}(v),
\]
the support function of the stated Cartesian product. Taking its
volume and dividing by $(|A||B|)^{r+s-1}$ proves multiplicativity.

Under $U$, a facet's area-normal vector $a_F\nu_F$ becomes
$|\det U|U^{-t}(a_F\nu_F)$. One way to see the area factor is to map
a prism of unit height over $F$: its volume multiplies by
$|\det U|$, while its height over the image facet becomes
$1/\|U^{-t}\nu_F\|$. Applying the facet formula to the transformed
segments proves the covariance identity. The determinant of the
linear map $|\det U|U^{-t}$ has absolute value
$|\det U|^{n-1}$, exactly the scale factor of $|UK|^{n-1}$.
Projection volumes are translation invariant.
\end{proof}

\begin{proposition}[Simplex value]\label{prop:proj-simplex}
For every positive integer $d$, any nondegenerate $d$-simplex satisfies
$R_d(T_d)=s_d$.
\end{proposition}

\begin{proof}
By affine invariance it suffices to use
$T_d=\operatorname{conv}(0,e_1,\ldots,e_d)$, of volume $1/d!$.
For $d\ge2$ its coordinate facets have area-normal vectors
$-e_i/(d-1)!$, and the opposite facet has vector
$w/(d-1)!$, where $w=e_1+\cdots+e_d$. The latter identity follows
either by its height and volume or by projecting that facet onto a
coordinate hyperplane. The facet formula shows that $\Pi T_d$ is a
translate of
\[
 \frac{[0,1]^d+[0,w]}{(d-1)!}.
\]
Let $Q=[0,1]^d$. Each fiber of $Q$ parallel to $w$ is an interval;
adding the segment $[0,w]$ extends its length by $\|w\|$ and leaves
the base projection unchanged. Fubini gives
\[
 |Q+[0,w]|=|Q|+\|w\|\,|\operatorname{proj}_{w^\perp}Q|_{d-1}
         =1+h_{\Pi Q}(w)=d+1.
\]
The last equality follows directly from the $2d$ unit-area facets of
the cube. Thus
\[
 R_d(T_d)=\frac{(d+1)(d!)^{d-1}}{((d-1)!)^d}
          =\frac{(d+1)d^d}{d!}=s_d.
\]
For $d=1$, $R_1(I)=2=s_1$ by the stated convention.
\end{proof}

Every simplex in a fixed dimension is an invertible affine image of
the standard one. Applying its affine map independently in each
factor, followed by a final ambient affine map, proves that the
geometric maximization over our class is precisely the finite
optimization \eqref{eq:proj-Pn}.

\section{Concavity and exact comparisons}

\begin{lemma}[Ratio monotonicity]\label{lem:proj-concavity}
The sequence $f_d=\log s_d$ is concave for $d\ge1$.
Its consecutive slopes satisfy
\[
 q_d:=\frac{s_{d+1}}{s_d}
   =\frac{d+2}{d+1}\left(1+\frac1d\right)^d,
 \qquad q_1=q_2=3>q_3>q_4>\cdots.
\]
\end{lemma}

\begin{proof}
Only monotonicity needs proof. For real $x\ge3$, let
\[
 h(x)=\log\frac{x+2}{x+1}
          +x\log\left(1+\frac1x\right).
\]
Direct differentiation gives
\[
 h'(x)=\log\left(1+\frac1x\right)
                   +\frac1{x+2}-\frac2{x+1},
 \qquad
 h''(x)=\frac{x^2-x-4}{x(x+1)^2(x+2)^2}>0.
\]
Since $h'(x)\to0$ as $x\to\infty$, strict increase of $h'$ implies
$h'(x)<0$ throughout $[3,\infty)$. Therefore $q_d$ strictly decreases
for integer $d\ge3$. Finally $q_1=q_2=3$ and $q_3=80/27<3$.
\end{proof}

The following comparisons are all between explicit positive rational
numbers:
\begin{align}
 Q_-&=\frac{s_{13}^{12}}{s_{12}^{13}}>1,
 &Q_+&=\frac{s_{13}^{14}}{s_{14}^{13}}>1,
 \label{eq:proj-peak-checks}\\
 G&=\frac{s_{14}^{8}}{s_{13}^{4}s_{12}^{5}}>1,
 &H&=\frac{s_{13}^{6}s_{12}^{4}}{s_{14}^{9}}>1.
 \label{eq:proj-branch-checks}
\end{align}
For a compact exact certificate, the ratios lie strictly between the
following rational endpoints. Every displayed integer is divided by
$10^6$.
\begin{center}
\begin{tabular}{lrr}
\toprule
Ratio & Lower numerator & Upper numerator\\
\midrule
$Q_-$ & 1019009 & 1019010\\
$Q_+$ & 1010131 & 1010132\\
$G$   & 1001039 & 1001040\\
$H$   & 1001185 & 1001186\\
$s_{12}^2/(s_{11}s_{13})$ & 1002550 & 1002551\\
$s_{14}^2/(s_{13}s_{15})$ & 1001959 & 1001960\\
$s_{13}^2/(s_{12}s_{14})$ & 1002226 & 1002227\\
\bottomrule
\end{tabular}
\end{center}
These statements can be checked by integer multiplication alone. For
example, the four main ratios have the prime-power forms
\begin{align*}
 Q_-&=\frac{5^2 7^{13}11\,13^{131}}{2^{290}3^{151}},
 &Q_+&=\frac{13^{181}}{2^{165}3^{18}5^{15}7^{156}11},\\
 G&=\frac{5^{10}7^{101}11}{2^{10}3^{47}13^{61}},
 &H&=\frac{3^{34}13^{85}}{2^{25}5^{11}7^{112}11}.
\end{align*}
The accompanying standard-library Python verifier stores each numerator
and denominator and both positive cross-multiplication differences. No
floating-point decision occurs in these certificates.

\begin{lemma}[The unique rate maximum]\label{lem:proj-peak}
The sequence $f_d/d$ strictly increases up to $d=13$ and strictly
decreases thereafter. Its piecewise-linear interpolation in the
numerator has the same property: if $f(x)$ is linear on each
$[d,d+1]$ and equals $f_d$ at integers, then $f(x)/x$ strictly increases
on $[1,13]$ and strictly decreases on $[13,\infty)$.
\end{lemma}

\begin{proof}
Set $a_d=d\log q_d-f_d$. Then
\[
 \frac{f_{d+1}}{d+1}-\frac{f_d}{d}
    =\frac{a_d}{d(d+1)},\qquad
 a_{d+1}-a_d=(d+1)(\log q_{d+1}-\log q_d)\le0.
\]
The two exact comparisons \eqref{eq:proj-peak-checks} say precisely that
$a_{12}>0$ and $a_{13}<0$. Monotonicity therefore gives the asserted
signs in all dimensions. On $(d,d+1)$, direct differentiation gives
$(f(x)/x)'=a_d/x^2$, proving the interpolated statement as well.
\end{proof}

\begin{proof}[Proof of Theorem~\ref{thm:proj-two}]
For fixed $k$, concavity gives
\[
 \sum_{i=1}^k f_{d_i}\le k f(n/k),
\]
and the balanced dimensions $a$ and $a+1$ attain equality. Thus the
maximum at that factor count is $B(n,k)$.
Since $k f(n/k)=n[f(n/k)/(n/k)]$, Lemma~\ref{lem:proj-peak} shows
that this expression increases with $k$ up to the real value $n/13$
and decreases thereafter. Only its two adjacent admissible integers
can maximize it. For $n<13$, the expression decreases throughout the
admissible factor counts, so $k=1$ is optimal. For $13\le n\le19$, the
remaining comparisons are the seven explicit rational comparisons
between one simplex and two balanced simplices. They give strict
preference to one simplex; at $n=20$ the two ten-dimensional factors
win, with the exact ratio established in \cite{OpenAIProjection}.
For completeness, the table below gives rational enclosures for the
two-factor value divided by the simplex value. All endpoints are
divided by $10^6$; the lower inequality is weak and the upper inequality
is strict.
\begin{center}
\begin{tabular}{crrr}
\toprule
$n$ & Factor dimensions & Lower numerator & Upper numerator\\
\midrule
13 & $(6,7)$ & 870778 & 870779\\
14 & $(7,7)$ & 893750 & 893751\\
15 & $(7,8)$ & 913688 & 913689\\
16 & $(8,8)$ & 935695 & 935696\\
17 & $(8,9)$ & 955049 & 955050\\
18 & $(9,9)$ & 976160 & 976161\\
19 & $(9,10)$ & 994936 & 994937\\
20 & $(10,10)$ & 1015230 & 1015231\\
\bottomrule
\end{tabular}
\end{center}
The verifier supplies the exact numerator and denominator for each
entry. At $n=14$ the lower endpoint is exact: the ratio is $143/160$.
\end{proof}

\begin{corollary}[Prescribed factor and facet counts]
\label{cor:proj-facets-budget}
Among simplex products of total dimension $n$ with exactly $k$ factors,
the maximum is $B(n,k)$. Such a product has exactly $n+k$ facets.
Consequently, for each $F$ with $n+1\le F\le2n$, the maximum among
these products with exactly $F$ facets is $B(n,F-n)$.
\end{corollary}

\begin{proof}
The fixed-factor claim is the concavity argument in the preceding
proof. A $d_i$-simplex has $d_i+1$ facets, and the facet classification
in Proposition~\ref{prop:proj-product} shows that the counts add under
Cartesian products. Their sum is $\sum_i(d_i+1)=n+k$.
\end{proof}

\section{The residue argument}

For $d\ge1$, define the nonnegative rate deficit
\[
 \delta_d=dL-f_d.
\]
In particular $\delta_{13}=0$, $\delta_{12}=\alpha$, and
$\delta_{14}=\beta$.

\begin{lemma}[Linear lower bounds for the deficits]
\label{lem:proj-deficit}
For positive integer $d$,
\[
 \delta_d\ge
 \begin{cases}
 (13-d)\alpha,&d\le12,\\
 0,&d=13,\\
 (d-13)\beta,&d\ge14.
 \end{cases}
\]
Equality is possible only at $d\in\{12,13,14\}$.
\end{lemma}

\begin{proof}
For $d\le12$, decreasing consecutive slopes give
\[
 f_{13}-f_d=\sum_{j=d}^{12}\log q_j
                  \ge(13-d)\log q_{12}.
\]
Since $f_{13}=13L$ and $\log q_{12}-L=\alpha$, this is the
first bound. It is strict if $d\le11$, because the sum then includes
a slope strictly greater than $\log q_{12}$. The argument for
$d\ge14$ uses
$f_d-f_{13}\le(d-13)\log q_{13}$ and
$L-\log q_{13}=\beta$, with strict inequality at $d\ge15$.
\end{proof}

For any factorization of $n$, put
\[
 A=\sum_i(d_i-13)_+,
 \qquad B=\sum_i(13-d_i)_+.
\]
If there are $k$ factors, then $n=13k+A-B$. Thus
$A-B\equiv r\pmod {13}$. Lemma~\ref{lem:proj-deficit} yields
\begin{equation}\label{eq:proj-residue-bound}
 nL-\sum_i f_{d_i}\ge A\beta+B\alpha
       \ge D_r:=\min\{r\beta,(13-r)\alpha\},
\end{equation}
where $D_0=0$. To prove the second inequality, remove
$\min(A,B)$ from both $A$ and $B$, which preserves the residue and
weakly lowers the positive cost. The remaining nonzero integer is
either congruent to $r$ or to $-r$ modulo $13$. Its least possible
positive value is therefore $r$ or $13-r$, respectively; an extra
multiple of $13$ only increases the cost.

The exact comparisons \eqref{eq:proj-branch-checks} are equivalent to
\[
 8\beta<5\alpha,\qquad 4\alpha<9\beta.
\]
They place the unique crossing between $r=8$ and $r=9$. Hence
\begin{equation}\label{eq:proj-Dr}
 D_r=\begin{cases}
 r\beta,&0\le r\le8,\\
 (13-r)\alpha,&9\le r\le12.
 \end{cases}
\end{equation}

\begin{proof}[Proof of Theorem~\ref{thm:proj-optimum}]
The rate assertion follows from Lemma~\ref{lem:proj-peak}.
Equation~\eqref{eq:proj-residue-bound} bounds every product by
$\rho^n e^{-D_r}$. If $r\le8$ and $n\ge14r$, the product of $r$
fourteen-dimensional simplices and $(n-14r)/13$ thirteen-dimensional
simplices attains that bound. If $r\ge9$ and
$n\ge12(13-r)$, the product with $13-r$ twelve-dimensional factors
and $(n-12(13-r))/13$ thirteen-dimensional factors attains it.
This proves the claimed formulas and their individual thresholds.

All equality conditions are strict away from the displayed choices.
The first inequality in \eqref{eq:proj-residue-bound} forces every
dimension to belong to $\{12,13,14\}$. Cancellation is strict when
$A,B>0$, so twelve- and fourteen-dimensional factors cannot coexist.
Additional multiples of $13$ in $A$ or $B$ are strictly more expensive.
Finally the two residue branches never tie. This proves uniqueness of
the dimension multiset.

For $1\le r\le8$, the last nonnegative dimension in that residue class
which fails $n\ge14r$ is $14r-13$; its largest value is $99$, at
$r=8$. The residue $r=0$ has no excluded positive dimension. For
$r\ge9$, the largest excluded dimension is at most $35$.
Thus every $n\ge100$ satisfies the relevant threshold. At $n=99$,
equality in the bound would require eight fourteen-dimensional
factors, whose total dimension already exceeds $99$. Strictness of
the equality conditions proves failure of the periodic formula there,
so the uniform threshold cannot be lowered.
\end{proof}

\section{Isolation and stability}

\begin{theorem}[Sharp isolation of the maximizing products]
\label{thm:proj-gap}
Let $n\ge100$. If a product of simplices $K$ of total dimension $n$
has a dimension multiset different from the maximizing multiset in
Theorem~\ref{thm:proj-optimum}, then
\[
 R_n(K)\le\frac{P_n}{G},\qquad
 G=\frac{s_{14}^{8}}{s_{13}^{4}s_{12}^{5}}
      \simeq1.00103981380358.
\]
No larger dimension-independent factor can replace $G$. For every
$n\equiv8\pmod {13}$ with $n\ge112$, equality is attained by
\[
 T_{12}^{\,5}\times T_{13}^{\,(n-60)/13}.
\]
\end{theorem}

\begin{proof}
Write $g=\log G=5\alpha-8\beta$ and
$h=\log H=9\beta-4\alpha$. The exact certificate table gives
\begin{equation}\label{eq:proj-gap-comparisons}
 \min\left\{h,\ \log\frac{s_{12}^2}{s_{11}s_{13}},\
 \log\frac{s_{14}^2}{s_{13}s_{15}},\ \alpha+\beta,
 \ 13\alpha,\ 13\beta\right\}>g>0.
\end{equation}
For a dimension $d\le11$, the remainder beyond the first bound in
Lemma~\ref{lem:proj-deficit} is
\[
 \delta_d-(13-d)\alpha
 =\sum_{j=d}^{12}(\log q_j-\log q_{12})
 \ge\log q_{11}-\log q_{12}
 =\log\frac{s_{12}^2}{s_{11}s_{13}}>g.
\]
For $d\ge15$, the analogous remainder is at least
$\log q_{13}-\log q_{14}
=\log(s_{14}^2/(s_{13}s_{15}))>g$.
Consequently any factor outside $\{12,13,14\}$ already produces a
loss strictly greater than $g$ beyond the universal residue bound.

It remains to consider products using only those three dimensions.
Their deficit is exactly $A\beta+B\alpha$. If both $A$ and $B$ are
positive, cancelling one from each reduces the cost by
$\alpha+\beta>g$, while the remaining cost is at least $D_r$.
If only one is positive but exceeds the least representative of its
residue, the extra cost is at least $13\alpha$ or $13\beta$, also
greater than $g$. The only remaining nonoptimal possibility is the
other residue branch. The difference between the two branch costs is
\[
 |r\beta-(13-r)\alpha|
      =|r(\alpha+\beta)-13\alpha|.
\]
As a function of integral $r$, its minimum occurs at one of the two
integers adjacent to the crossing, namely $8$ and $9$. The values
there are $g$ and $h$, respectively, and $h>g$. Thus every nonoptimal
product loses at least $g$ in logarithmic value. Exponentiating proves
the inequality.

For $r=8$, the optimal branch uses eight fourteen-dimensional factors.
The other branch uses five twelve-dimensional factors, giving the
displayed product whenever $n\ge112$. Its logarithmic deficit exceeds
the optimum by exactly $5\alpha-8\beta=g$. This proves sharpness.
\end{proof}

\begin{corollary}[Bounded exceptions in near maximizers]
\label{cor:proj-stability}
For $n\ge100$ and a simplex product with factor dimensions $d_i$,
put $\varepsilon=\log(P_n/R_n(K))\ge0$. Then
\[
 \sum_i|d_i-13|\le
 \frac{\varepsilon+D_r}{\min\{\alpha,\beta\}}.
\]
In particular, within any fixed multiplicative factor of the maximum,
the number of factors of dimension different from $13$ is bounded
independently of the total dimension.
\end{corollary}

\begin{proof}
The left side is $A+B$ and
$\varepsilon+D_r=\sum_i\delta_{d_i}
\ge A\beta+B\alpha
\ge\min\{\alpha,\beta\}(A+B)$.
\end{proof}

\section{Periodicity and asymptotics}

The preceding theorem implies the exact recurrence
\begin{equation}\label{eq:proj-rec}
 P_{n+13}=s_{13}P_n\qquad(n\ge100),
\end{equation}
and hence a rational ordinary generating function:
\[
 \sum_{n\ge1}P_n z^n
 =\sum_{n=1}^{99}P_nz^n
   +\frac{\sum_{n=100}^{112}P_nz^n}{1-s_{13}z^{13}}.
\]
Here the coefficients are rational optimal values; no counting
interpretation is asserted.

For $n\ge100$ with residue $r$, the excess over the simplex has the
exact expression
\begin{equation}\label{eq:proj-asymp-exact}
 \frac{P_n}{s_n}
 =e^{-D_r}\rho^n\frac{n!}{(n+1)n^n}.
\end{equation}
Stirling's expansion therefore gives, uniformly over the thirteen
residues,
\begin{equation}\label{eq:proj-asymp}
 \frac{P_n}{s_n}
 =e^{-D_r}\sqrt{\frac{2\pi}{n}}
       \left(\frac\rho e\right)^n
 \left(1-\frac{11}{12n}+\frac{265}{288n^2}
              -\frac{47839}{51840n^3}+O(n^{-4})\right).
\end{equation}
In particular the optimal exponential improvement within this entire
product class is
$\rho/e\simeq1.03372825553$; the subexponential prefactor is an
explicit period-$13$ constant times $n^{-1/2}$.
Stirling's formula also gives $s_d^{1/d}\to e$. The strict decrease
of these rates after dimension $13$ therefore proves $\rho>e$
without relying on the displayed decimal approximation.

More generally, for each fixed positive integer $J$, the standard
logarithmic Stirling expansion \cite{DLMFStirling} yields
\begin{align*}
 \log\frac{P_n}{s_n}
 ={}&n(\log\rho-1)-\frac12\log n
       +\frac12\log(2\pi)-D_r-\log(1+1/n)\\
 &+\sum_{j=1}^{J}
   \frac{B_{2j}}{2j(2j-1)n^{2j-1}}
   +O(n^{-2J-1}).
\end{align*}
This determines every algebraic correction explicitly. The finite
periodic term is exact, and the only asymptotic input is Stirling's
classical expansion.

\section{Validation and further questions}

The geometric identities in \eqref{eq:proj-inputs}, the dimension-twenty
counterexample, and the existence of an exponential excess are source
results. The contribution here is the optimization of the product
class: the unique block size $13$, the all-dimensional two-candidate
formula, the exact eventual residue classes and sharp threshold, the
sharp isolation factor, and the periodic asymptotic prefactor.
A targeted search did not locate these specific optimization statements;
that search is evidence of a plausible new refinement rather than a
proof of priority. None of the arguments assumes the source collection's
more ambitious claims about arbitrary convex bodies.

The exact arithmetic artifact verifies seven rational interval
certificates and independently computes $P_n$ for $1\le n\le300$ by
the recurrence
\[
 P_0=1,\qquad P_n=\max_{1\le d\le n}s_dP_{n-d}.
\]
Every dimension agrees with the two-candidate theorem, and every
attainable residue formula agrees as well. These finite computations
check constants and implementations; the proofs above establish the
claims for all dimensions.

A separate implementation retains the two largest distinct rational
values at every dimension and tracks their dimension multisets. It
checks all dimensions through $260$, finds a unique maximizing multiset
in each, and verifies the isolation bound from dimension $100$ onward.
Within that range its minimum observed ratio equals $G$ exactly, at
$n=112,125,138,151,164,177,190,203,216,229,242,255$.
For example, it recovers
\[
 \operatorname{Opt}(99)=(12,12,12,12,12,13,13,13)
\]
and the runner-up at dimension $112$ with dimensions
$(12,12,12,12,12,13,13,13,13)$.
Both programs use rational arithmetic for all decisions; decimals in
their reports are only displays. The full comparison outputs and
source-provenance record accompany the manuscript.

Several research directions follow directly.
\begin{enumerate}
 \item Can one exhibit a convex body $K$ in some dimension $d$ with
 $R_d(K)^{1/d}>\rho$? Such a body would improve the best exponential
 rate supplied by simplex products, through the same product identity.
 \item Does the supremum of $R_d(K)^{1/d}$ over all dimensions and all
 convex bodies occur at a finite dimension? The present theorem
 answers this restricted to simplex factors.
 \item How does the optimum change with a prescribed vertex budget
 or simultaneous vertex and facet budgets? The vertex count is
 $\prod_i(d_i+1)$, while Corollary~\ref{cor:proj-facets-budget} solves
 the facet-count constraint alone. The combined problem imposes an
 additional nonlinear restriction on the dimension partition.
 \item Can the deficit method be extended to factors drawn from
 several geometrically distinct families? A unique best rate would
 suggest eventual periodic corrections governed by a finite
 optimization over dimension residues, but concavity may fail.
 \item Does a geometric perturbation of an optimal product improve
 its normalized projection-body volume while preserving its
 combinatorial type? The sharp gap above concerns changes of factor
 dimensions, and imposes no obstruction to such perturbations.
 \item What are the exact second- and third-largest product values
 in every sufficiently large dimension? The sharp uniform gap locates
 one infinite family of runner-up products, but a full residue-by-residue
 classification of further values would describe the local discrete
 optimization landscape more precisely.
 \item For centrally symmetric factor families, is there an analogous
 unique optimal block size and finite residue rule? The relevant
 factor values must first be established or bounded, and they need
 not satisfy the concavity exploited here.
 \item Can the complete argument be formalized, including the facet
 formula and affine covariance, so that the arithmetic certificates
 and the geometric reduction are checked in one proof system?
 The present artifact verifies the finite arithmetic and supplies
 ordinary mathematical proofs of the claims in all dimensions;
 it does not contain a Lean formalization.
\end{enumerate}

\begin{thebibliography}{9}
\bibitem{OpenAIProjection}
OpenAI,
\emph{A product counterexample to the simplex maximum for
projection-body volume}, 24 September 2026.
\href{https://github.com/openai/math/tree/adc7f1241b42e322a6451854ab7e4b4c146bf78a/preprints/A-product-counterexample-to-the-simplex-maximum-for-projection-body-volume-September-24-2026}{Manuscript and source files in the openai/math repository}.
The inspected \texttt{build/main.tex} has Git blob
\texttt{c32187535cddb5c38f28d332997dc0578cbad586}.

\bibitem{Brannen1996}
N. S. Brannen,
\emph{Volumes of projection bodies},
Mathematika \textbf{43} (1996), no.~2, 255--264.
\href{https://doi.org/10.1112/S002557930001175X}{doi:10.1112/S002557930001175X}.

\bibitem{FengReverse}
Y. Feng, S. Hu, W. Liu, and L. Xu,
\emph{On the Reverse Projection Inequality}, preprint, 2026.
\href{https://doi.org/10.5281/zenodo.22037130}{doi:10.5281/zenodo.22037130}.
Primary repository record published 21 August 2026.

\bibitem{DLMFStirling}
NIST Digital Library of Mathematical Functions,
\emph{Section 5.11: Asymptotic Expansions}, especially
equations 5.11.1--5.11.4 and Section 5.11(ii).
\url{https://dlmf.nist.gov/5.11}.
Accessed 8 October 2026.
\end{thebibliography}
\end{document}
